\chapter{复现与证据索引}
\label{app:repro}

\section{代码仓库与可重现产物}
综合代码仓库为 \RepositoryLink。仓库根目录的 README 提供依赖安装、
工作目录、课程数据目录和入口命令。LaTeX 主文件位于 \code{report/main.tex}，
图表脚本位于 \code{report/scripts/make_figures.py}。本地正式版以
\code{cover.local.tex} 提供完整封面字段。公开检出可以使用同一主文件
编译去除个人资料的版本。

\begin{table}[H]
\centering\small
\caption{复查各任务的代码与结果入口}
\begin{tabularx}{\linewidth}{@{}P{21mm}YY@{}}
\toprule
任务 & 主要代码 & 可核查的记录 \\
\midrule
幼苗 & \code{task1/plantcls.py}，\code{train.py}，\code{predict.py}
 & \code{report/data/plant/} 中配置、折外汇总与逐轮日志 \\
表情 & \code{task2/src/ds.py}，\code{train.py}，\code{blend_search.py}
 & \code{report/data/fer/} 中代表性日志、融合权重和消融 \\
定位 & \code{task3/solution/src/model.py}，\code{infer.py}，\code{vascx_fovea.py}
 & \code{report/data/fovea/} 中运行配置和逐折记录 \\
分割 & \code{task4/drive_vessel/}，\code{inference.py}，\code{BEST_CONFIG}
 & 实验日志、原生训练日志及冻结提交哈希 \\
\bottomrule
\end{tabularx}
\end{table}

已归档的 \code{report/data/kaggle_scores.csv} 保存各赛题完成状态的
提交名称、编号、时间和两侧分数。哈希列对能与本地具名文件匹配的产物
保存 SHA-256，供教师核对本地提交文件。历史上名称重复的
\code{submission.csv} 仅凭同名不能确认文件内容。主表对关键分数保留
提交编号，在 \code{report/evidence_audit.md} 记录证据路径和计算口径。

\section{复现顺序}
\begin{enumerate}
\item 使用课程 Kaggle 页面或教师发放的文件获取四项数据。保持训练标签、测试
图像和提交样例的目录结构。使用项目 README 指定的工作目录和相对数据路径。
\item 首先用发布的脚本检查输入与提交格式。再按需要运行代表性训练、推理和
融合。历史权重由本地材料或经许可的上游渠道获取。
\item 对任意已归档提交用哈希校验其字节，对 CSV 检查字段与样例顺序。
分割提交增加 RLE 解码，中心凹提交核查坐标轴、单位与图像边界。
\item 在仓库根目录生成图表，在报告目录编译主文件。构建自动运行 BibTeX。
\begin{quote}\small
\texttt{python report/scripts/make\_figures.py}\\
\texttt{cd report}\\
\texttt{latexmk -xelatex -outdir=build main.tex}
\end{quote}
\end{enumerate}

\section{复查时容易混淆的口径}
\begin{itemize}
\item 本地分类报告中的 OOF 表示单张训练图由其对应的验证折模型预测。
测试时每张图可由所有折共同预测，两者的投票成员数不同。
\item 中心凹 RMSE 汇总按样本数量加权平方误差。直接平均各折 RMSE
会改变聚合方式。距离分档的样本内得分与嵌套验证结果各自记录。
\item DRIVE 官方权重的 epoch 40 表示上游 checkpoint。原生 U-Net
运行的实际轮次也为 40，其验证最佳在第 15 轮，两条记录来自不同模型。
\item 参考论文建立方法来源。课程数据量、折划分、训练轮次、后处理和
分数均由项目或平台记录核查。
\end{itemize}

\subsection{环境与外部资源}
交付核查在 Windows 的 Python 3.13 环境执行，PyTorch 为
2.14.0+cu126，torchvision 为 0.29.1+cu126，timm 为 1.0.30，
OpenCV 为 5.0.0，albumentations 为 2.0.8。该环境用于检查
发布入口、提交格式和已存档的分割概率。原任务三训练记录报告的 GPU
为 RTX 4070 Laptop、显存 8 GB，且保存了 PyTorch、timm、
OpenCV 与增强库版本。其他运行的硬件与依赖仅在各自记录可证实时列出。

\begin{table}[H]
\centering\small
\caption{复现中需要另行准备的输入}
\begin{tabularx}{\linewidth}{@{}P{23mm}YP{44mm}@{}}
\toprule
任务 & 本地数据入口 & 其他依赖 \\
\midrule
幼苗 & \code{task1/dataset-for-task2/} & timm 预训练骨干 \\
表情 & \code{task2/2026task2/} 或 \code{FER_PROJECT_ROOT} & timm 骨干；验证和测试概率缓存 \\
定位 & \code{task3/2026task3/detection/} 或 \code{FOVEA_DATA_DIR} & 血管模型、VascX 与其预处理组件 \\
分割 & \code{task4/2026task-4/segmentation/} & 上游 FR-UNet checkpoint \\
\bottomrule
\end{tabularx}
\end{table}
